\clearpage
\sectionkicker{Source-backed technical notes}
\subsection*{Paper Watch — DiffusionGemma: Technical Notes}
この欄は記事本文で圧縮した一次資料上の情報を、比較・再検証しやすい形へ展開したものである。新しい外部情報は追加せず、Selection済みEvidenceのchronology、normalized claim、limitations、source URLのみを再配置する。
\medskip
\subsection*{Theme at a glance}
\addcontentsline{toc}{subsection}{Theme at a glance}
\begin{center}
\footnotesize
\begin{tabularx}{\linewidth}{@{}p{0.20\linewidth}p{0.10\linewidth}p{0.14\linewidth}X@{}}
\toprule
Artifact & Role & Type & Objective chronology \\
\midrule
DiffusionGemma & PRIMARY & OPEN\_WEIGHT & 2026-07-31T16:11:46Z (PAPER\_PUBLICATION) \\
\bottomrule
\end{tabularx}
\end{center}
\normalsize

\begin{technicalnote}{DiffusionGemma}{PRIMARY}
\begin{tabularx}{\linewidth}{@{}>{\bfseries}p{0.22\linewidth}X@{}}
Organization & - \\
Artifact type & OPEN\_WEIGHT \\
Chronology & 2026-07-31T16:11:46Z (PAPER\_PUBLICATION) \\
\end{tabularx}
\smallskip
{\bfseries 一次資料から整理したtechnical points}
\begin{itemize}[leftmargin=1.5em,itemsep=0.35em]
\item \textbf{Author claim}: The paper describes DiffusionGemma as an experimental open-weight discrete-diffusion language model derived from a Gemma-family MoE model and using blockwise iterative refinement.
\item \textbf{Author claim}: The authors report speed/quality results for their sampler, including high single-H100 decoding throughput; these numbers are study results rather than independent reproduction.
\end{itemize}
{\bfseries 読む際の境界}
\begin{itemize}[leftmargin=1.5em,itemsep=0.35em]
\item \textbf{分析上の留意点}: The technical report is an author-controlled evaluation of an experimental model; speed and quality results depend on its sampler, model configuration, and hardware setup.
\end{itemize}
{\bfseries Primary source}
\begin{itemize}[leftmargin=1.5em,itemsep=0.25em]
\item \url{http://arxiv.org/abs/2608.00146v1}
\end{itemize}
{\scriptsize\color{SurveyMuted}Source-bound record: \texttt{evidence:SP-2026-M07:item-arxiv-2608-00146v1-7b3dbb80d8}.}
\end{technicalnote}
\medskip
